\section{Discussion}
\label{sec:discussion}

\subsection{Quasi-static family of states and branching initiation}
\label{subsec:discussion-quasistatic}

In the present formulation, the V-shaped interfacial shear crack is a
prescribed stationary configuration: the length $l$ of each crack branch
remains fixed, whereas the equilibrium process-zone length $d_i$ changes
with the outer-loading parameter $C$. The solution therefore defines the
one-parameter quasi-static family
\begin{equation*}
  C\longrightarrow d_i(C)
  \longrightarrow
  \bigl\{\delta_i(C),W_i(C),J_i(C)\bigr\}.
\end{equation*}
Here, quasi-static refers to a family of equilibrium zone configurations
in which inertia is neglected.

This family can be used directly to determine the critical load for
branching initiation. If a critical opening $\delta_{i\mathrm{c}}$ or a
critical zone-related energy parameter $J_{i\mathrm{c}}$ is specified
independently for the material, the critical state is defined by
\begin{equation*}
  \delta_i(C_c)=\delta_{i\mathrm{c}}
  \qquad\text{or}\qquad
  J_i(C_c)=J_{i\mathrm{c}}.
\end{equation*}
Once the corresponding critical material characteristic is known, these
relations determine $C_c$. Subsequent propagation and the finite path of
the branched crack require a separate formulation.

\subsection{Competing fracture mechanisms}
\label{subsec:discussion-competing}

Two distinct local mechanisms are possible for the stationary V-shaped
crack. The first is associated with its corner $O$: the singular corner
field supports a tensile process zone in one of the materials, and
attainment of a critical $\delta_i$ or $J_i$ marks branching initiation.
The second mechanism is localized at the outer tips of the V-shaped
crack. For those tips, Nazarenko and Kipnis
\cite{NazarenkoKipnis2022Equilibrium} determined the Mode~II stress
intensity factor $K_{\mathrm{II}}$ and formulated the interfacial
fracture condition
\begin{equation*}
  K_{\mathrm{II}}=K_{\mathrm{IIc}}.
\end{equation*}

Branching from the corner and interfacial propagation from the outer tips
should therefore be regarded as competing local mechanisms of the same
stationary geometry, rather than as prescribed sequential stages. Their
respective critical loads are governed by different material
characteristics: Mode~II interfacial fracture resistance at the outer tips
and the critical process-zone characteristic, $\delta_i$ or $J_i$, at
the corner.

The small-scale assumption $d_i\ll l$ permits the two local regions to be
treated separately at leading order. The field of the stationary V-shaped
crack serves as the outer field for the process-zone problem near $O$, and
the feedback of the small zone on the fields at the outer crack tips is
neglected. A coupled boundary-value problem would be required for a
quantitative analysis of the competition when $d_i/l$ is no longer small.

\subsection{Energy and deformation criteria}
\label{subsec:discussion-criteria}

According to Eq.~\eqref{eq:energy-opening-relation}, for fixed geometry
and elastic properties,
\begin{equation*}
  J_i=\chi_i\sigma_i\delta_i.
\end{equation*}
Thus, $J_i$ and $\delta_i$ are not independent measures of the current
zone state: both are uniquely determined by the same equilibrium length
$d_i$. The deformation criterion
$\delta_i=\delta_{i\mathrm{c}}$ and the energy criterion
$J_i=J_{i\mathrm{c}}$ therefore select the same critical state when
\begin{equation*}
  J_{i\mathrm{c}}
  =\chi_i\sigma_i\delta_{i\mathrm{c}}.
\end{equation*}

This equivalence is specific to the present one-parameter model and fixed
geometry. It does not imply a universal identity between two independently
measured material characteristics for arbitrary fracture problems, because
$\chi_i$ depends on the corner geometry and elastic contrast. In
applications, either critical characteristic may be used, provided it is
established independently and the two criteria are not treated as
independent constraints on the same model state.

\subsection{Residual corner field and local energy flux}
\label{subsec:discussion-flux}

In the configuration containing the process zone, the field at the corner
$O$ remains singular. If $H_i$ is the amplitude of the leading corner
eigenfield, then
\begin{equation*}
  \sigma\sim H_i r^{\lambda_i}.
\end{equation*}
The corresponding elastic strain-energy density is of order
\begin{equation*}
  w\sim\frac{H_i^2}{E_i}r^{2\lambda_i},
\end{equation*}
and the local configurational energy flux through a contour of
characteristic radius $R$ scales as
\begin{equation}
\label{eq:local-elastic-flux-scaling}
  J_i^{(\mathrm{el})}(R)
  \sim\frac{H_i^2}{E_i}R^{2\lambda_i+1}.
\end{equation}

For the physically admissible regimes identified in the numerical
analysis, $-1/2<\lambda_i<0$. Hence $2\lambda_i+1>0$ and
\begin{equation*}
  \lim_{R\to0}J_i^{(\mathrm{el})}(R)=0.
\end{equation*}
The residual stress singularity at $O$ therefore does not generate a
finite nonzero local elastic contribution to the zone energy parameter.
In this sense, the corner singularity remaining after zone formation is
weaker than the classical $r^{-1/2}$ crack singularity.

This distinction also identifies the limit of the additive representation
discussed by Kaminsky and Dudyk~\cite{KaminskyDudyk2026Criteria}. For an
existing crack, the outer field has a singularity exponent with real part
$-1/2$, and a finite elastic energy flux can enter the overall energy
balance together with the process-zone contribution. For the present
corner problem, $\lambda_i>-1/2$ implies that
$J_i^{(\mathrm{el})}(R)\to0$, so no separate finite elastic term is
required in the branching-initiation criterion. The zone-related parameter
$J_i=\dd W_i/\dd d_i$ can therefore be used without such an additive
correction; this does not make $J_i$ an independently established global
contour integral.

\subsection{Range of applicability and generalization}
\label{subsec:discussion-scope}

The analysis is based on the scale separation $d_i\ll l\ll L$. We use
$d_i/l<0.1$ as the principal range for quantitative application. Portions
of the numerical curves for which $d_i/l$ exceeds $0.1$ and reaches
approximately $0.18$ near the maxima remain useful for showing trends, but
they represent extrapolations of the local approximation. In those
regimes, higher-order terms of the outer asymptotic field and interaction
between the local regions near the corner and the outer crack tips become
increasingly important.

The local character of the solution also facilitates extension to a
broader class of problems. For a developed symmetric V-shaped interfacial
crack with finite branches, the factor $CQ_i$ that supplies the outer-field
amplitude in the local problem can be replaced by the amplitude, or
generalized stress intensity factor, of the corresponding corner singular
field. Such an extension remains consistent provided that the process zone
is small and the same local eigenproblem applies. The model can therefore
be used to analyze branching initiation at a kink in an interface. A
related application is splitting of a solid by a sharp wedge when its
local boundary-value problem reduces to an analogous corner configuration.
